\paragraph{теорема Вейєрштрасса.} Фундаментальною основою для використання поліномів у задачах наближення функцій є класичне твердження, яке гарантує принципову можливість такого наближення. Воно стверджує, що будь-яка неперервна функція на обмеженому інтервалі може бути наближена поліномами з довільною наперед заданою точністю. Цей факт формалізується у вигляді відомої теореми\cite{trefethen2013approximation}.

\begin{statement}[Апроксимаційна теорема Вейєрштрасса]
	\label{st:weierstrass}
	Нехай $f$~--- неперервна функція на відрізку $[-1, 1]$, і нехай задано довільне як завгодно мале додатне число $\varepsilon > 0$. Тоді існує такий поліном $p$, для якого виконується умова:
	\begin{equation}\label{eq:weierstrass_condition}
		\|f - p\| < \varepsilon.
	\end{equation}
\end{statement}

Важливим наслідком теореми Вейєрштрасса є те, що вона встановлює принципову можливість поліноміального наближення навіть для вкрай негладких функцій. До цієї категорії належать, наприклад, такі функції як $x \sin(x^{-1})$ або навіть $\sin(x^{-1})\sin(1/\sin(x^{-1}))$ ~\cite{trefethen2013approximation}.